\subsection{Drawing}\label{sec:drawing-task}

Despite being trained on only textual data, LMs have been shown to be able to structure their representations of perceptual concepts such as size and color~\citepia{abdou-etal-2021-language, patel2022mapping, zhang-etal-2020-language-embeddings, ilharco-etal-2021-probing}  in a way that credibly mirrors the physical world. Recent LMs can even generate plausible drawings of objects using code such  as TikZ and SVG~\cite{bubeck2023sparks,zhang2023controllable}. We evaluate the visual understanding of LMs by asking them to generate code for drawing various objects in the Processing language, which \citet{sharma2024vision} found the LMs to be more adept in. Psychological studies have shown that humans have the ability to rotate mental representations of objects~\cite{shepard1971mental, vandenberg1978mental}. For the counterfactual settings, we similarly ask the LM to generate code that draws the same object, but rotated or vertically flipped. We disallow the use of functions such as \texttt{rotate} to prevent  shortcut solutions  (see~\S\ref{sec:overestimation} for further discussion).
As with the spatial reasoning task (\S\ref{sec:spatial-task}), an ideal model should be robust to these settings.
For the CCC, we ask the model to draw a straight  line at the top of the canvas in addition to the object; a flipped/rotated line thus signifies an understanding of the transformations.
